% from ega4-2-fr.tex
\oldpage[IV-2]{81}

\begin{proposition}[4.8.3]
\label{IV.2.4.8.3-fr}
Les applications canoniques (4.8.2.1) à (4.8.2.5) sont injectives.
\end{proposition}

\begin{proof}
On sait en effet que $p_X$ est fidèlement plat et quasi-compact; le fait que
(4.8.2.5) soit injectif résulte donc de ce que $p_X$ est surjectif.
L'injectivité de (4.8.2.1) résulte de (2.2.2), celle de (4.8.2.2) de (2.2.7),
celle de (4.8.2.3) de (2.2.15) et enfin celle de (4.8.2.4) de (2.2.16).
\end{proof}

\begin{definition}[4.8.4]
\label{IV.2.4.8.4-fr}
Les notations étant celles de (4.8.2), on dit qu'un
sous-$\mathcal O_{X_{(K)}}$-Module quasi-cohérent de
$\mathscr F\otimes_kK$ (resp. un homomorphisme
$\mathscr F\otimes_kK\to\mathscr G\otimes_kK$, resp. un sous-préschéma de
$X_{(K)}$, resp. un $K$-morphisme $X_{(K)}\to Y_{(K)}$, resp. un
sous-ensemble de $X_{(K)}$) est \emph{défini sur $K'$} s'il appartient à
l'image de l'application (4.8.2.1) (resp. (4.8.2.2), resp. (4.8.2.3), resp.
(4.8.2.4), resp. (4.8.2.5)). On dit alors que $K'$ est un \emph{corps de
définition} de l'objet considéré.
\end{definition}

Il est clair que $K$ lui-même est corps de définition d'un élément de l'un
quelconque des cinq ensembles d'arrivée des applications (4.8.2.1) à
(4.8.2.5). On notera toutefois que pour un $K$-préschéma $Z$, dire par exemple
qu'un $\mathcal O_Z$-Module quasi-cohérent $\mathscr H$ est «~défini sur un
sous-corps $K'$ de $K$~» n'a de sens que si $Z$ \emph{a été donné} sous la
forme $X_{(K)}$ pour un préschéma $X$ sur un sous-corps $k$ de $K$, et si
$\mathscr H$ est un sous-$\mathcal O_Z$-Module d'un $\mathcal O_Z$-Module qui
\emph{a été donné} sous la forme $\mathscr F\otimes_kK$, où $\mathscr F$ est
un $\mathcal O_X$-Module quasi-cohérent; on a des remarques analogues pour les
autres notions. Il résulte de (4.8.2.6) que si un élément de l'un des ensembles
d'arrivée des applications de (4.8.2) est défini sur $K'$, il est aussi défini
sur tout corps $K''$ tel que $K'\subset K''\subset K$. Enfin, avec les
notations de (4.8.2.6), si $K_1$ est une extension de $K$, pour qu'un élément
de l'ensemble d'arrivée de $\varepsilon_{K,K'}$ soit défini sur $K'$, il
faut et il suffit que son image par $\varepsilon_{K_1,K}$ soit définie sur
$K'$, en vertu de la relation
$\varepsilon_{K_1,K'}=\varepsilon_{K_1,K}\circ\varepsilon_{K,K'}$.

\begin{proposition}[4.8.5]
\label{IV.2.4.8.5-fr}
Avec les notations de (4.8.2), soit $(X_\alpha)$ un recouvrement ouvert de
$X$. Pour qu'un élément $\mathscr H\in\Phi(\mathscr F\otimes_kK)$ (resp.
$u\in\operatorname{Hom}(\mathscr F\otimes_kK,\mathscr G\otimes_kK)$, resp.
$Z\in\mathfrak S(X_{(K)})$, resp. $M\in\mathfrak P(X_{(K)})$) soit défini sur
$K'$, il faut et il suffit que pour tout $\alpha$,
$\mathscr H|(X_\alpha)_{(K)}$ (resp. $u|(X_\alpha)_{(K)}$,
$Z\cap(X_\alpha)_{(K)}$, $M\cap(X_\alpha)_{(K)}$) soit défini sur $K'$.
\end{proposition}

\begin{proof}
Cela résulte aussitôt de l'injectivité des applications (4.8.2.1), (4.8.2.2),
(4.8.2.3) et (4.8.2.5): par exemple, si pour tout $\alpha$, il y a un
sous-$\mathcal O_{(X_\alpha)_{(K')}}$-Module quasi-cohérent
$\mathscr H'_\alpha$ de $(\mathscr F|X_\alpha)\otimes_kK'$ tel que
$\mathscr H|(X_\alpha)_{(K)}=\mathscr H'_\alpha\otimes_{K'}K$, pour deux
indices $\alpha$, $\beta$ quelconques, on aura
\[
\bigl(\mathscr H'_\alpha|(X_\alpha\cap X_\beta)_{(K')}\bigr)
 \otimes_{K'}K
=\bigl(\mathscr H'_\beta|(X_\alpha\cap X_\beta)_{(K')}\bigr)
 \otimes_{K'}K,
\]
donc nécessairement
$\mathscr H'_\alpha|(X_\alpha\cap X_\beta)_{(K')}
=\mathscr H'_\beta|(X_\alpha\cap X_\beta)_{(K')}$ quels que soient
$\alpha$ et $\beta$ et il y a par suite un sous-$\mathcal O_{X_{(K')}}$-Module
quasi-cohérent $\mathscr H'$ de $\mathscr F\otimes_kK'$ tel que
$\mathscr H'|(X_\alpha)_{(K')}=\mathscr H'_\alpha$ pour tout $\alpha$, donc
$\mathscr H=\mathscr H'\otimes_{K'}K$. On raisonne de même dans les autres
cas.
\end{proof}

\begin{lemma}[4.8.6]
\label{IV.2.4.8.6-fr}
Soient $k$ un corps, $K$ une extension de $k$, $A$ une $k$-algèbre, $M$ un
$A$-module, $\overline N$ un sous-$A_{(K)}$-module de $M_{(K)}$, $K'$ une
sous-extension de $K$. Les conditions suivantes sont équivalentes:
\begin{enumerate}
\item[a)] $\overline N$ est de la forme $N'\otimes_{K'}K$, où $N'$ est un
sous-$A_{(K')}$-module de $M_{(K')}$. 
\item[b)] Si $X=\operatorname{Spec}(A)$, le
$\mathcal O_{X_{(K)}}$-Module quasi-cohérent $\widetilde{\overline N}$ sur
$X_{(K)}=\operatorname{Spec}(A_{(K)})$ est défini sur $K'$.
\item[c)] $\overline N$ est de la forme $N'\otimes_{K'}K$, où $N'$ est un
sous-$K'$-espace vectoriel de $M_{(K')}$. 
\end{enumerate}
\end{lemma}

